\section{H}

% halten (to hold, stop)
\begin{verbentry}{
  style = {default},
  toc = {halten (to hold, stop)},
  name = {halten},
  english = {to hold, stop},
  type = {irregular},
  auxiliary = {haben}
}
 
    
     \vspace{5pt}

    \hrule
    
    \vspace{5pt}

  \begin{tabular}{rl}
     \multicolumn{2}{l}{\textbf{PRÄSENS} }\\
    ich & \textbf{halte}\\
    du & \textbf{hältst}\\
    er/sie/es & \textbf{hält}\\
    wir & \textbf{halten}\\
    ihr & \textbf{haltet}\\
    sie/Sie & \textbf{halten}\\
  \end{tabular}
  \hfill
     \begin{tabular}{rl}
    \multicolumn{2}{l}{\textbf{PRÄTERITUM} }\\
    ich & \textbf{hielt}\\
    du & \textbf{hieltst}\\
    er/sie/es & \textbf{hielt}\\
    wir & \textbf{hielten}\\
    ihr & \textbf{hieltet}\\
    sie/Sie & \textbf{hielten}\\
  \end{tabular}
  \hfill
  \begin{tabular}{rl}
     \multicolumn{2}{l}{\textbf{IMPERATIV} }\\
   & \textbf{halt(e) (du)!}\\
   & \textbf{halten wir!}\\
   & \textbf{haltet (ihr)!}\\
   & \textbf{halten Sie!}\\
   & \vspace{10pt}
  \end{tabular}

    \vspace{10pt}

 \begin{tabular}{rl}
     \multicolumn{2}{l}{\textbf{PERFEKT}}\\
  &  haben + \textbf{gehalten}\\
  \end{tabular}
\hfill
   \begin{tabular}{rl}
   
    \multicolumn{2}{l}{\textbf{FUTURE I} }\\
   & werden + \textbf{halten}\\
 
  \end{tabular}
\hfill
   \begin{tabular}{rl}
      \multicolumn{2}{l}{\textbf{PLUSQUAMPERFEKT}}\\
   & hatten + \textbf{gehalten}\\
  \end{tabular}
  
  \vspace{10pt}



\begin{tabular}{lll}
\multicolumn{3}{l}{\textbf{MIT PRÄPOSITIONEN}}\\
& \textbf{---} & \\
\end{tabular}
\hfill
\begin{tabular}{lll}
\multicolumn{3}{l}{\textbf{TRENNBARE VERBEN}}\\
& \textbf{---} & \\
\end{tabular}

  \vspace{-5pt}
\end{verbentry}

% heben (to lift / raise)
\begin{verbentry}{
  style = {default},
  toc = {heben (to lift / raise)},
  name = {heben},
  english = {to lift / raise},
  type = {irregular},
  auxiliary = {haben}
}
 
    
     \vspace{5pt}

    \hrule
    
    \vspace{5pt}

  \begin{tabular}{rl}
     \multicolumn{2}{l}{\textbf{PRÄSENS} }\\
    ich & \textbf{hebe}\\
    du & \textbf{hebst}\\
    er/sie/es & \textbf{hebt}\\
    wir & \textbf{heben}\\
    ihr & \textbf{hebt}\\
    sie/Sie & \textbf{heben}\\
  \end{tabular}
  \hfill
     \begin{tabular}{rl}
    \multicolumn{2}{l}{\textbf{PRÄTERITUM} }\\
    ich & \textbf{hob}\\
    du & \textbf{hobst}\\
    er/sie/es & \textbf{hob}\\
    wir & \textbf{hoben}\\
    ihr & \textbf{hobt}\\
    sie/Sie & \textbf{hoben}\\
  \end{tabular}
  \hfill
  \begin{tabular}{rl}
     \multicolumn{2}{l}{\textbf{IMPERATIV} }\\
   & \textbf{heb(e) (du)!}\\
   & \textbf{heben wir!}\\
   & \textbf{hebt (ihr)!}\\
   & \textbf{heben Sie!}\\
   & \vspace{10pt}
  \end{tabular}

    \vspace{10pt}

 \begin{tabular}{rl}
     \multicolumn{2}{l}{\textbf{PERFEKT}}\\
   & haben + \textbf{gehoben}\\
  \end{tabular}
\hfill
   \begin{tabular}{rl}
   
    \multicolumn{2}{l}{\textbf{FUTURE I} }\\
   & werden + \textbf{heben}\\
 
  \end{tabular}
\hfill
   \begin{tabular}{rl}
      \multicolumn{2}{l}{\textbf{PLUSQUAMPERFEKT}}\\
   & hatten + \textbf{gehoben}\\
  \end{tabular}

   \vspace{10pt}

   \begin{tabular}{lll}
      \multicolumn{3}{l}{\textbf{MIT PRÄPOSITIONEN}}\\
& \textbf{heben auf} + Akk & (to pick up / keep / store)\\
& \textbf{heben von} + Dat & (to lift from)\\
& \textbf{heben aus} + Dat & (to lift out of)\\
\vspace{10pt}
   \end{tabular}
  \hfill
   \begin{tabular}{lll}
      \multicolumn{3}{l}{\textbf{TRENNBARE VERBEN}}\\
 & \textbf{ab$|$heben} & (to lift off / withdraw money)\\
 & \textbf{auf$|$heben} & (to pick up / cancel / keep)\\
 & \textbf{an$|$heben} & (to begin / raise slightly)\\
 & \textbf{heraus$|$heben} & (to lift out)\\
  \end{tabular}

\vspace{-5pt}
\end{verbentry}

% heiraten (to marry)
\begin{verbentry}{
  style = {default},
  toc = {heiraten (to marry)},
  name = {heiraten},
  english = {to marry},
  type = {regular},
  auxiliary = {haben}
}
 

\vspace{5pt}
\hrule
\vspace{5pt}

\begin{tabular}{rl}
\multicolumn{2}{l}{\textbf{PRÄSENS}}\\
ich & \textbf{heirate}\\
du & \textbf{heiratest}\\
er/sie/es & \textbf{heiratet}\\
wir & \textbf{heiraten}\\
ihr & \textbf{heiratet}\\
sie/Sie & \textbf{heiraten}\\
\end{tabular}
\hfill
\begin{tabular}{rl}
\multicolumn{2}{l}{\textbf{PRÄTERITUM}}\\
ich & \textbf{heiratete}\\
du & \textbf{heiratetest}\\
er/sie/es & \textbf{heiratete}\\
wir & \textbf{heirateten}\\
ihr & \textbf{heiratetet}\\
sie/Sie & \textbf{heirateten}\\
\end{tabular}
\hfill
\begin{tabular}{rl}
\multicolumn{2}{l}{\textbf{IMPERATIV}}\\
& \textbf{heirate (du)!}\\
& \textbf{heiraten wir!}\\
& \textbf{heiratet (ihr)!}\\
& \textbf{heiraten Sie!}\\
\end{tabular}

\vspace{10pt}

\begin{tabular}{rl}
\multicolumn{2}{l}{\textbf{PERFEKT}}\\
& haben + \textbf{geheiratet}\\
\end{tabular}
\hfill
\begin{tabular}{rl}
\multicolumn{2}{l}{\textbf{FUTURE I}}\\
& werden + \textbf{heiraten}\\
\end{tabular}
\hfill
\begin{tabular}{rl}
\multicolumn{2}{l}{\textbf{PLUSQUAMPERFEKT}}\\
& hatten + \textbf{geheiratet}\\
\end{tabular}

\vspace{10pt}

\begin{tabular}{lll}
\multicolumn{3}{l}{\textbf{MIT PRÄPOSITIONEN}}\\
& \textbf{heiraten} + Akk & (to marry someone)\\
& \textbf{heiraten mit} + Dat & (to marry, formal)\\
\end{tabular}
\hfill
\begin{tabular}{lll}
\multicolumn{3}{l}{\textbf{TRENNBARE VERBEN}}\\
& \textbf{---} & \\
\end{tabular}

\vspace{-5pt}
\end{verbentry}

% heißen (to be called)
\begin{verbentry}{
  style = {default},
  toc = {heißen (to be called)},
  name = {heißen},
  english = {to be called},
  type = {irregular},
  auxiliary = {haben}
}
 
    
     \vspace{5pt}

    \hrule
    
    \vspace{5pt}

  \begin{tabular}{rl}
     \multicolumn{2}{l}{\textbf{PRÄSENS} }\\
    ich & \textbf{heiße}\\
    du & \textbf{heißt}\\
    er/sie/es & \textbf{heißt}\\
    wir & \textbf{heißen}\\
    ihr & \textbf{heißt}\\
    sie/Sie & \textbf{heißen}\\
  \end{tabular}
  \hfill
     \begin{tabular}{rl}
    \multicolumn{2}{l}{\textbf{PRÄTERITUM} }\\
    ich & \textbf{hieß}\\
    du & \textbf{hießest}\\
    er/sie/es & \textbf{hieß}\\
    wir & \textbf{hießen}\\
    ihr & \textbf{hießt}\\
    sie/Sie & \textbf{hießen}\\
  \end{tabular}
  \hfill
  \begin{tabular}{rl}
     \multicolumn{2}{l}{\textbf{IMPERATIV} }\\
   & \textbf{heiß(e) (du)!}\\
   & \textbf{heißen wir!}\\
   & \textbf{heißt (ihr)!}\\
   & \textbf{heißen Sie!}\\
   & \vspace{10pt}
  \end{tabular}

    \vspace{10pt}

 \begin{tabular}{rl}
     \multicolumn{2}{l}{\textbf{PERFEKT}}\\
   & haben + \textbf{geheißen}\\
  \end{tabular}
\hfill
   \begin{tabular}{rl}
   
    \multicolumn{2}{l}{\textbf{FUTURE I} }\\
   & werden + \textbf{heißen}\\
 
  \end{tabular}
\hfill
   \begin{tabular}{rl}
      \multicolumn{2}{l}{\textbf{PLUSQUAMPERFEKT}}\\
   & hatten + \textbf{geheißen}\\
  \end{tabular}
  
  \vspace{10pt}



\begin{tabular}{lll}
\multicolumn{3}{l}{\textbf{MIT PRÄPOSITIONEN}}\\
& \textbf{---} & \\
\end{tabular}
\hfill
\begin{tabular}{lll}
\multicolumn{3}{l}{\textbf{TRENNBARE VERBEN}}\\
& \textbf{---} & \\
\end{tabular}

  \vspace{-5pt}
\end{verbentry}

% helfen (to help)
\begin{verbentry}{
  style = {dat},
  toc = {helfen (to help)},
  name = {helfen},
  english = {to help},
  type = {irregular, + Dativ},
  auxiliary = {haben}
}
 
    
     \vspace{5pt}

    \hrule
    
    \vspace{5pt}

  \begin{tabular}{rl}
     \multicolumn{2}{l}{\textbf{PRÄSENS} }\\
    ich & \textbf{helfe}\\
    du & \textbf{hilfst}\\
    er/sie/es & \textbf{hilft}\\
    wir & \textbf{helfen}\\
    ihr & \textbf{helft}\\
    sie/Sie & \textbf{helfen}\\
  \end{tabular}
  \hfill
     \begin{tabular}{rl}
    \multicolumn{2}{l}{\textbf{PRÄTERITUM} }\\
    ich & \textbf{half}\\
    du & \textbf{halfst}\\
    er/sie/es & \textbf{half}\\
    wir & \textbf{halfen}\\
    ihr & \textbf{halft}\\
    sie/Sie & \textbf{halfen}\\
  \end{tabular}
  \hfill
  \begin{tabular}{rl}
     \multicolumn{2}{l}{\textbf{IMPERATIV} }\\
   & \textbf{hilf (du)!}\\
   & \textbf{helfen wir!}\\
   & \textbf{helft (ihr)!}\\
   & \textbf{helfen Sie!}\\
   & \vspace{10pt}
  \end{tabular}

    \vspace{10pt}

 \begin{tabular}{rl}
     \multicolumn{2}{l}{\textbf{PERFEKT}}\\
  &  haben + \textbf{geholfen}\\
  \end{tabular}
\hfill
   \begin{tabular}{rl}
   
    \multicolumn{2}{l}{\textbf{FUTURE I} }\\
   & werden + \textbf{helfen}\\
 
  \end{tabular}
\hfill
   \begin{tabular}{rl}
      \multicolumn{2}{l}{\textbf{PLUSQUAMPERFEKT}}\\
   & hatten + \textbf{geholfen}\\
  \end{tabular}
  
  
\vspace{10pt}

\begin{tabular}{lll}
\multicolumn{3}{l}{\textbf{MIT PRÄPOSITIONEN}}\\
& \textbf{---} & \\
\end{tabular}
\hfill
\begin{tabular}{lll}
\multicolumn{3}{l}{\textbf{TRENNBARE VERBEN}}\\
& \textbf{---} & \\
\end{tabular}
\vspace{-5pt}
\end{verbentry}

% holen (to fetch)
\begin{verbentry}{
  style = {acc},
  toc = {holen (to fetch)},
  name = {holen},
  english = {to fetch},
  type = {regular, + Akkusativ},
  auxiliary = {haben}
}
 
    
     \vspace{5pt}

    \hrule
    
    \vspace{5pt}

  \begin{tabular}{rl}
     \multicolumn{2}{l}{\textbf{PRÄSENS} }\\
    ich & \textbf{hole}\\
    du & \textbf{holst}\\
    er/sie/es & \textbf{holt}\\
    wir & \textbf{holen}\\
    ihr & \textbf{holt}\\
    sie/Sie & \textbf{holen}\\
  \end{tabular}
  \hfill
     \begin{tabular}{rl}
    \multicolumn{2}{l}{\textbf{PRÄTERITUM} }\\
    ich & \textbf{gholte}\\
    du & \textbf{holtest}\\
    er/sie/es & \textbf{holte}\\
    wir & \textbf{holten}\\
    ihr & \textbf{holtet}\\
    sie/Sie & \textbf{holten}\\
  \end{tabular}
  \hfill
  \begin{tabular}{rl}
     \multicolumn{2}{l}{\textbf{IMPERATIV} }\\
   & \textbf{hol(e) (du)!}\\
   & \textbf{holen wir!}\\
   & \textbf{holt (ihr)!}\\
   & \textbf{holen Sie!}\\
   & \vspace{10pt}
  \end{tabular}

    \vspace{10pt}

 \begin{tabular}{rl}
     \multicolumn{2}{l}{\textbf{PERFEKT}}\\
   & haben + \textbf{geholt}\\
  \end{tabular}
\hfill
   \begin{tabular}{rl}
   
    \multicolumn{2}{l}{\textbf{FUTURE I} }\\
   & werden + \textbf{holen}\\
 
  \end{tabular}
\hfill
   \begin{tabular}{rl}
      \multicolumn{2}{l}{\textbf{PLUSQUAMPERFEKT}}\\
   & hatten + \textbf{geholt}\\
  \end{tabular}
  
  \vspace{10pt}

   \begin{tabular}{lll}
      \multicolumn{3}{l}{\textbf{MIT PRÄPOSITIONEN}}\\
& \textbf{holen für} + Akk & (to fetch something for someone)\\
& \textbf{holen aus} + Dat & (to take out of / fetch from inside)\\
&\textbf{holen nach} + Dat & (to fetch to a location)\\
&\textbf{abholen von} + Akk & (to pick up from)
  \end{tabular}
  \hfill
   \begin{tabular}{lll}
      \multicolumn{3}{l}{\textbf{TRENNBARE VERBEN}}\\
 & \textbf{ab$|$holen} & (to pick up / collect)\\
 & \textbf{auf$|$holen} & (to catch up)\\
     & \textbf{ein$|$holen} & (to catch up / retrieve)\\
     & \textbf{nach$|$holen} & (to make up / catch up later)
  \end{tabular}


  \vspace{-5pt}
\end{verbentry}

% hoffen (to hope)
\begin{verbentry}{
  style = {acc},
  toc = {hoffen (to hope)},
  name = {hoffen},
  english = {to hope},
  type = {regular, + Akkusativ},
  auxiliary = {haben}
}
 
    
     \vspace{5pt}

    \hrule
    
    \vspace{5pt}

  \begin{tabular}{rl}
     \multicolumn{2}{l}{\textbf{PRÄSENS} }\\
    ich & \textbf{hoffe}\\
    du & \textbf{hoffst}\\
    er/sie/es & \textbf{hofft}\\
    wir & \textbf{hoffen}\\
    ihr & \textbf{hofft}\\
    sie/Sie & \textbf{hoffen}\\
  \end{tabular}
  \hfill
     \begin{tabular}{rl}
    \multicolumn{2}{l}{\textbf{PRÄTERITUM} }\\
    ich & \textbf{hoffte}\\
    du & \textbf{hofftest}\\
    er/sie/es & \textbf{hoffte}\\
    wir & \textbf{hofften}\\
    ihr & \textbf{hofftet}\\
    sie/Sie & \textbf{hofften}\\
  \end{tabular}
  \hfill
  \begin{tabular}{rl}
     \multicolumn{2}{l}{\textbf{IMPERATIV} }\\
   & \textbf{hoff(e) (du)!}\\
   & \textbf{hoffen wir!}\\
   & \textbf{hofft (ihr)!}\\
   & \textbf{hoffen Sie!}\\
   & \vspace{10pt}
  \end{tabular}

    \vspace{10pt}

 \begin{tabular}{rl}
     \multicolumn{2}{l}{\textbf{PERFEKT}}\\
   & haben + \textbf{gehofft}\\
  \end{tabular}
\hfill
   \begin{tabular}{rl}
   
    \multicolumn{2}{l}{\textbf{FUTURE I} }\\
   & werden + \textbf{hoffen}\\
 
  \end{tabular}
\hfill
   \begin{tabular}{rl}
      \multicolumn{2}{l}{\textbf{PLUSQUAMPERFEKT}}\\
   & hatten + \textbf{gehofft}\\
  \end{tabular}
  
  
\vspace{10pt}

\begin{tabular}{lll}
\multicolumn{3}{l}{\textbf{MIT PRÄPOSITIONEN}}\\
& \textbf{---} & \\
\end{tabular}
\hfill
\begin{tabular}{lll}
\multicolumn{3}{l}{\textbf{TRENNBARE VERBEN}}\\
& \textbf{---} & \\
\end{tabular}
\vspace{-5pt}
\end{verbentry}

% hören (to hear)
\begin{verbentry}{
  style = {acc},
  toc = {hören (to hear)},
  name = {hören},
  english = {to hear},
  type = {regular, + Akkusativ},
  auxiliary = {haben}
}
 
    
     \vspace{5pt}

    \hrule
    
    \vspace{5pt}

  \begin{tabular}{rl}
     \multicolumn{2}{l}{\textbf{PRÄSENS} }\\
    ich & \textbf{höre}\\
    du & \textbf{hörst}\\
    er/sie/es & \textbf{hört}\\
    wir & \textbf{hören}\\
    ihr & \textbf{hört}\\
    sie/Sie & \textbf{hören}\\
  \end{tabular}
  \hfill
     \begin{tabular}{rl}
    \multicolumn{2}{l}{\textbf{PRÄTERITUM} }\\
    ich & \textbf{hörte}\\
    du & \textbf{hörtest}\\
    er/sie/es & \textbf{hörte}\\
    wir & \textbf{hörten}\\
    ihr & \textbf{hörtet}\\
    sie/Sie & \textbf{hörten}\\
  \end{tabular}
  \hfill
  \begin{tabular}{rl}
     \multicolumn{2}{l}{\textbf{IMPERATIV} }\\
   & \textbf{hör(e) (du)!}\\
   & \textbf{hören wir!}\\
   & \textbf{hört (ihr)!}\\
   & \textbf{hören Sie!}\\
   & \vspace{10pt}
  \end{tabular}

    \vspace{10pt}

 \begin{tabular}{rl}
     \multicolumn{2}{l}{\textbf{PERFEKT}}\\
    &haben + \textbf{gehört}\\
  \end{tabular}
\hfill
   \begin{tabular}{rl}
   
    \multicolumn{2}{l}{\textbf{FUTURE I} }\\
    &werden + \textbf{hören}\\
 
  \end{tabular}
\hfill
   \begin{tabular}{rl}
      \multicolumn{2}{l}{\textbf{PLUSQUAMPERFEKT}}\\
    &hatten + \textbf{gehört}\\
  \end{tabular}
  

  \vspace{10pt}

   \begin{tabular}{lll}
      \multicolumn{3}{l}{\textbf{MIT PRÄPOSITIONEN}}\\
& \textbf{hören auf} + Akk & (to listen to / obey / pay attention to)\\
& \textbf{hören von} + Dat & (to hear about / receive news)\\
\vspace{30pt}
  \end{tabular}
  \hfill
   \begin{tabular}{lll}
      \multicolumn{3}{l}{\textbf{TRENNBARE VERBEN}}\\
 & \textbf{ab$|$hören} & (to listen to / examine)\\
 & \textbf{an$|$hören} & (to hear completely)\\
 & \textbf{auf$|$hören} & (to stop / quit)\\
  & \textbf{mit$|$hören} & (to listen in / eavesdrop)\\
  & \textbf{nach$|$hören} & (to check / listen again)\\
   & \textbf{zu$|$hören} & (to listen carefully)\\
  \end{tabular}

  \vspace{-5pt}
\end{verbentry}

% husten (to cough)
\begin{verbentry}{
  style = {default},
  toc = {husten (to cough)},
  name = {husten},
  english = {to cough},
  type = {regular},
  auxiliary = {haben}
}


\vspace{5pt}
\hrule
\vspace{5pt}

\begin{tabular}{rl}
\multicolumn{2}{l}{\textbf{PRÄSENS}}\\
ich & \textbf{huste}\\
du & \textbf{hustest}\\
er/sie/es & \textbf{hustet}\\
wir & \textbf{husten}\\
ihr & \textbf{hustet}\\
sie/Sie & \textbf{husten}\\
\end{tabular}
\hfill
\begin{tabular}{rl}
\multicolumn{2}{l}{\textbf{PRÄTERITUM}}\\
ich & \textbf{hustete}\\
du & \textbf{hustetest}\\
er/sie/es & \textbf{hustete}\\
wir & \textbf{husteten}\\
ihr & \textbf{hustetet}\\
sie/Sie & \textbf{husteten}\\
\end{tabular}
\hfill
\begin{tabular}{rl}
\multicolumn{2}{l}{\textbf{IMPERATIV}}\\
& \textbf{huste (du)!}\\
& \textbf{husten wir!}\\
& \textbf{hustet (ihr)!}\\
& \textbf{husten Sie!}\\
\end{tabular}

\vspace{10pt}

\begin{tabular}{rl}
\multicolumn{2}{l}{\textbf{PERFEKT}}\\
& haben + \textbf{gehustet}\\
\end{tabular}
\hfill
\begin{tabular}{rl}
\multicolumn{2}{l}{\textbf{FUTURE I}}\\
& werden + \textbf{husten}\\
\end{tabular}
\hfill
\begin{tabular}{rl}
\multicolumn{2}{l}{\textbf{PLUSQUAMPERFEKT}}\\
& hatten + \textbf{gehustet}\\
\end{tabular}

\vspace{10pt}

\begin{tabular}{lll}
\multicolumn{3}{l}{\textbf{MIT PRÄPOSITIONEN}}\\
& \textbf{husten in} + Akk & (to cough into)\\
& \textbf{husten vor} + Dat & (to cough from / because of)\\
\end{tabular}
\hfill
\begin{tabular}{lll}
\multicolumn{3}{l}{\textbf{TRENNBARE VERBEN}}\\
& \textbf{aus$|$husten} & (to cough out)\\
\end{tabular}
\vspace{-5pt}
\end{verbentry}
